% Evaluation, Conclusions, and Future Work (15 marks)
% - Does it clearly summarise outcomes;
% - How was the quality of choice of evaluation approach(es), including choice of metrics, data
% collection and preparation;
% - Does it provide a critical assessment of the outcomes with reference to the original
% objectives and literature reviewed;
% - Does it identify of any limitation and appropriate recommendations;
% - Is there indication of continuity and/or extension in the form of future work.
\chapter{Evaluation}
\label{chp:evaluation}
% \paragraph{Evaluation setup}
% We are going to evaluate the three approaches on the 56 Solomon datasets. For each
% instance, we will generate dynamic instances using multiple random seeds and benchmark
% them by running:
% For each instance, we will generate dynamic variants by using three different seeds, three cutoff times and
% three dod. instances will be generated beforehand and the models will be evaluated on the same instance. For a total of 3024 instances.
% Oracle has a 300s (5 minutes) budget and it solves the static variant. ORTools has 30 seconds to begin with and 5 seconds for re-opt both in hybrid and pure contexts.
% AI is with greedy decoding and augmentation. 8 augmentations.
% 
% \begin{itemize}
%     \item the heuristics (OR-Tools and PyVRP) with short time budgets;
%     \item AI models (RouteFinder and all the variants developed);
%     \item hybrid methods combining the above approaches.
% \end{itemize}
{\color{red} Use Ablation table in evaluation for training variants}

\section{Reproducibility and Implementation Details}

To support transparent comparison and independent replication, this section
summarises the computational environment, benchmark protocol, solver
configurations, randomness controls, and software versions used throughout the
study.

\subsection{Computational Environment}

All experiments were executed on the hardware platforms reported in
Table~\ref{tab:hardware}. Runtime comparisons should be interpreted with respect
to the reported machine specifications.

\begin{table}[H]
\centering
\caption{Hardware environment used for training and evaluation.}
\label{tab:hardware}
\begin{tabular}{llllll}
\toprule
Machine & CPU & Cores & RAM & GPU / VRAM & Operating System \\
\midrule
Training Node & [CPU MODEL] & [N] & [RAM] & [GPU MODEL] / [VRAM] & [OS] \\
Evaluation Node & [CPU MODEL] & [N] & [RAM] & [GPU MODEL] / [VRAM] & [OS] \\
Oracle Node & [CPU MODEL] & [N] & [RAM] & [GPU MODEL] / [VRAM] & [OS] \\
\bottomrule
\end{tabular}
\end{table}

\subsection{Runtime Budgets}

All compared methods were evaluated under common time-budget constraints.

\begin{table}[H]
\centering
\caption{Runtime budget policy.}
\label{tab:budgets}
\begin{tabular}{lll}
\toprule
Setting & Budget & Notes \\
\midrule
Static benchmark & [X seconds] & Per full instance solve \\
Dynamic benchmark & [Y seconds] & Per re-optimisation event \\
Oracle reference & [Z seconds] & Full-information offline solve \\
Hybrid methods & Shared budget & AI inference + refinement included \\
\bottomrule
\end{tabular}
\end{table}

Unless otherwise stated, one-off costs such as model loading, file parsing, and
library initialisation were excluded. Warm-start conversion overheads were
[counted / excluded], and this policy was applied consistently across all hybrid
runs.

\subsection{Dataset Protocol}

The benchmark was based on Solomon VRPTW instances. Table~\ref{tab:dataset}
should list the exact instance families and subsets used.

\begin{table}[H]
\centering
\caption{Dataset protocol.}
\label{tab:dataset}
\begin{tabular}{llll}
\toprule
Category & Instances Used & Customer Variants & Purpose \\
\midrule
Static evaluation & [LIST] & 50 / 75 / 100 & Benchmarking \\
Dynamic evaluation & [LIST] & 50 / 75 / 100 & Online benchmarking \\
AI training (general) & Generated & 50 / 75 & Training \\
AI training (Solomon-shaped) & [LIST / Generated] & 50 / 75 & Training \\
Validation & [LIST / Generated] & 50 / 75 & Model selection \\
\bottomrule
\end{tabular}
\end{table}

If Solomon-shaped generators were derived from evaluation instances, this should
be disclosed and justified. If separate training and test families were used,
the partition should be explicitly listed.

\subsection{Randomness and Seed Control}

To reduce variance and ensure fair comparison, randomness was controlled through
fixed seeds.

\begin{table}[H]
\centering
\caption{Random seed policy.}
\label{tab:seeds}
\begin{tabular}{ll}
\toprule
Component & Seed(s) \\
\midrule
Dynamic scenario generation & [SEED LIST] \\
OR-Tools solver randomness & [SEED / NONE] \\
PyVRP randomness & [SEED LIST] \\
RouteFinder training & [SEED LIST] \\
Inference augmentation / sampling & [SEED / DETERMINISTIC] \\
Scenario replications per instance/DoD & [N] \\
\bottomrule
\end{tabular}
\end{table}

Where stochastic methods were used, repeated runs should be aggregated through
mean, median, or paired comparisons.

\subsection{Benchmark Hyperparameters}

Table~\ref{tab:hyperparams} summarises global benchmark parameters.

\begin{table}[H]
\centering
\caption{Benchmark hyperparameters.}
\label{tab:hyperparams}
\begin{tabular}{ll}
\toprule
Parameter & Value \\
\midrule
Degree of Dynamism levels & $\{[...]\}$ \\
cutoff\_ratio ($\alpha$) & [VALUE] \\
Reveal-time distribution & Uniform $[0,\alpha T]$ \\
Soft lateness penalty (OR-Tools) & [VALUE] \\
Soft lateness penalty (RouteFinder) & [VALUE] \\
Hybrid shared-time policy & [DESCRIPTION] \\
Default inference augmentation & [VALUE] \\
\bottomrule
\end{tabular}
\end{table}

\subsection{Solver Configurations}

\subsubsection*{OR-Tools}

\begin{table}[H]
\centering
\caption{OR-Tools configuration.}
\label{tab:ortools}
\begin{tabular}{ll}
\toprule
Parameter & Value \\
\midrule
Version & [VERSION] \\
First solution strategy & [e.g. PATH\_CHEAPEST\_ARC] \\
Local search metaheuristic & [e.g. GUIDED\_LOCAL\_SEARCH] \\
Neighbourhood operators & [LIST] \\
Time limit & [VALUE] \\
Seed & [VALUE / NONE] \\
Soft TW penalty & [VALUE] \\
Capacity penalty (if any) & [VALUE / NONE] \\
\bottomrule
\end{tabular}
\end{table}

\subsubsection*{PyVRP}

\begin{table}[H]
\centering
\caption{PyVRP configuration.}
\label{tab:pyvrp}
\begin{tabular}{ll}
\toprule
Parameter & Value \\
\midrule
Version & [VERSION] \\
Population size & [VALUE] \\
Stopping criterion & [TIME / ITERATIONS] \\
Repair / education settings & [VALUE] \\
Seed & [VALUE] \\
Time cap & [VALUE] \\
Penalty management & [DESCRIPTION] \\
\bottomrule
\end{tabular}
\end{table}

\subsubsection*{Hybrid Method}

\begin{table}[H]
\centering
\caption{Hybrid warm-start configuration.}
\label{tab:hybrid}
\begin{tabular}{ll}
\toprule
Parameter & Value \\
\midrule
AI proposer & RouteFinder \\
Refinement solver & OR-Tools / PyVRP \\
Warm-start representation & \texttt{Solution} object \\
Translation overhead counted & [YES / NO] \\
Fallback on infeasible seed & [YES / NO] \\
Shared time budget & [YES] \\
\bottomrule
\end{tabular}
\end{table}

\subsection{RouteFinder Training Configuration}

\begin{table}[H]
\centering
\caption{RouteFinder training parameters.}
\label{tab:rftrain}
\begin{tabular}{ll}
\toprule
Parameter & Value \\
\midrule
Architecture version & [VERSION / COMMIT] \\
Epochs by variant & [LIST] \\
Batch size & [VALUE] \\
Instances per epoch & [VALUE] \\
Validation size & [VALUE] \\
Optimiser & Adam \\
Learning rate & [VALUE] \\
Weight decay & [VALUE] \\
Checkpoint frequency & [VALUE] \\
Model selection rule & Final / Best validation \\
Training seeds & [VALUE] \\
\bottomrule
\end{tabular}
\end{table}

\subsection{Software Stack}

\begin{table}[H]
\centering
\caption{Software versions.}
\label{tab:software}
\begin{tabular}{ll}
\toprule
Component & Version \\
\midrule
Python & [VERSION] \\
PyTorch & [VERSION] \\
CUDA & [VERSION] \\
OR-Tools & [VERSION] \\
PyVRP & [VERSION] \\
RouteFinder & [COMMIT / RELEASE] \\
NumPy / Pandas & [VERSION] \\
Operating system & [VERSION] \\
\bottomrule
\end{tabular}
\end{table}

\subsection{Stopping Criteria}

All optimisation procedures were terminated using explicit stopping rules.

\begin{itemize}
    \item \textbf{OR-Tools:} fixed wall-clock time limit of [X] seconds.
    \item \textbf{PyVRP:} terminated after [X] seconds or [Y] iterations.
    \item \textbf{Hybrid methods:} terminated when shared budget was exhausted.
    \item \textbf{RouteFinder training:} terminated after scheduled epochs.
\end{itemize}

\subsection{Statistical Analysis}

Performance comparisons were conducted across repeated scenarios and benchmark
instances. Unless otherwise stated, reported results should include means and
standard deviations, or medians and interquartile ranges when distributions are
non-normal.

For paired multi-method comparisons, non-parametric procedures such as the
Wilcoxon signed-rank test (two methods) or Friedman test with post-hoc Nemenyi
comparisons (multiple methods) are recommended. Confidence intervals should be
reported where feasible.

\subsection{Code Availability}

The full benchmarking framework, solver wrappers, training modifications, and
configuration files should be released in a public repository. The repository
should include:

\begin{itemize}
    \item exact seeds;
    \item solver parameter files;
    \item environment definitions;
    \item training scripts;
    \item raw results;
    \item plotting and statistical-analysis scripts.
\end{itemize}
